\documentclass[11pt]{article}
\usepackage{amsmath,amssymb,amsthm,hyperref,geometry}\geometry{margin=1in}
\newtheorem{theorem}{Theorem}\newtheorem{proposition}{Proposition}
\title{Separated grades and an admissible changing-shape interaction}
\date{15 November 2026}\begin{document}\maketitle
\paragraph{Status.} Scoped general construction with written proofs. The first theorem is categorical, independent of kernels or a particular language. A Set-diagram instance then supplies nontrivial changing-shape maps under a precise preservation criterion. This is a separated-grade family, not arbitrary interfering content or a new independent O1 criterion. Standard indexed monoidal and Kan-extension techniques are not claimed as novel.

\section{General indexed monoidal input}
Let $\mathcal C$ be a pure category, $K:\mathcal C^{\mathrm{op}}\to\mathsf{MonCat}$ and $L:\mathcal C\to\mathsf{MonCat}$ coherent pseudofunctors with strong monoidal transition functors and monoidal composition/unit comparison natural isomorphisms. Here $\mathsf{MonCat}$ has monoidal natural transformations as its 2-cells. Symmetry is not required. Write $u^*$ and $v_!$ for their actions. Take the context Grothendieck fibration with extension equal to its projection, and the content Grothendieck opfibration with extension equal to its projection. Thus $e_AP=A$ and $d_BQ=B$: the grades are retained native objects but do not themselves change the underlying pure interface. Unit and accumulation are the respective monoidal units and products.

The literal modules are
\[
 M(A,B)=K_A\times\mathcal C(A,B),\quad N(A,B)=\mathcal C(A,B)\times L_B,
 \quad H(A,B)=K_A\times\mathcal C(A,B)\times L_B,
\]
with hom sets of $\mathcal C$ regarded as discrete categories. A native mixed cell at fixed endpoints and pure body $f$ is an arbitrary pair of grade arrows; changing endpoints uses $P\to u^*P'$ and $v_!Q\to Q'$ with $vf=f'u$.

\begin{theorem}[Separated-grade native interaction]
These presentations admit a full directed crossing $N\odot M\Rightarrow M\odot N$. The induced composition on the literal mixed family is
\[
 (P,f,Q);(R,g,S)
   =\bigl(P\otimes f^*R,\ gf,\ g_!Q\otimes S\bigr),\tag{S}
\]
with unit $(I_{K_A},\mathrm{id}_A,I_{L_A})$. All native arrows and pure endpoint squares are retained.
\end{theorem}
\begin{proof}
For a source representative with middle $Y$, written $(f:A\to Y,Q\in L_Y);(P\in K_Y,g:Y\to B)$, define its crossed literal object as $(f^*P,gf,g_!Q)$. A middle balance along $u:Y\to Y'$ compares $f^*u^*P'$ with $(uf)^*P'$ and $(gu)_!Q$ with $g_!u_!Q$, using exactly the pseudofunctor composition comparisons. The same comparisons and transition functors act on all native grade maps, so naturality and balance identity/composition hold. Retained comma normalization converts the output to $M\odot N$ without identifying native maps.

Crossing a unit grade uses the strong monoidal unit comparisons. Crossing an accumulated context grade uses $f^*(P\otimes R)\cong f^*P\otimes f^*R$; crossing accumulated content uses $g_!(Q\otimes S)\cong g_!Q\otimes g_!S$. These are the four required comparisons. Their diagrams reduce to the monoidal coherence of $K,L$, the strong monoidal transition axioms, and pseudofunctor coherence. More explicitly every iterated word has context normal form consisting of its context grades in execution order, each pulled to the initial interface along its preceding pure composite; its content normal form consists of content grades in execution order, each pushed to the final interface along its subsequent pure composite. Every route uses the same grade order and transition composite. No symmetry or interchange of two grades within one doctrine is used. Strong monoidal coherence gives the same map to these parenthesized normal forms, including unit insertions, and proves all pseudodistributive diagrams. This also verifies (S) on objects and arrows; its associator has exactly those two normal forms.
\end{proof}

No additional mixed Beck--Chevalley or projection formula is invoked in this family: pullback acts only on $K$ grades, pushforward only on $L$ grades, and both extensions equal the pure projection. The theorem would not justify omitting those conditions when a grade changes its extension interface or carries shared effects.

The input hypothesis is substantive: merely lax monoidal transition maps would not yield the stated invertible comparisons. Nor does this theorem identify the two grade doctrines or claim that they describe a hidden joint physical system.

\section{A changing-shape pure category}
Let $\mathcal P$ have small categories as objects and those functors $r:J\to J'$ for which the usual left Kan extension
\[
 \operatorname{Lan}_r:\mathbf{Set}^J\to\mathbf{Set}^{J'}
\]
preserves finite products, including terminal objects, as its arrows. All these Kan extensions exist in the declared universes. Identities belong to $\mathcal P$, and it is closed under composition: the canonical comparison $\operatorname{Lan}_s\operatorname{Lan}_r\cong\operatorname{Lan}_{sr}$ transfers finite-product preservation. Equip these functors with the canonical strong cartesian monoidal comparisons. Their coherence follows from product universality and the canonical Kan composition comparisons.

Set $K_J=L_J=\mathbf{Set}^J$ with pointwise cartesian product. The context action is restriction $r^*$, which preserves products. The content action is $\operatorname{Lan}_r$, which preserves products by admission. These define the inputs to the theorem and hence a genuine separated-grade changing-shape native interaction. Its mixed objects retain a context diagram, an admitted shape functor, and a content diagram. This is not obtained by assuming every functor between shapes admissible.

\begin{proposition}[Checkable admission and a nontrivial fold]
If every incoming comma $(r\downarrow j')$ has Set-colimits commuting with finite products (is sifted), then $r$ is admitted. For a fold $J\to1$, admission is exactly siftedness of $J$. In particular the fold from the increasing chain of natural numbers is admitted, while the fold from two discrete objects is not.
\end{proposition}
\begin{proof}
Pointwise Kan extension is the colimit over the incoming comma of the restricted diagram. Restriction preserves pointwise finite products, so the stated property of its index implies preservation by the Kan extension. For a terminal target this is exactly the defining finite-product-colimit property of siftedness.

For the natural-number chain, a pair of colimit classes has representatives at stages $n,m$, which can both be transported to $\max(n,m)$, proving surjectivity of the comparison from the colimit of products. Equality of two pairs of classes is witnessed at some later stages in each coordinate; a common later stage witnesses equality of the pair, proving injectivity. The same argument handles any finite number of coordinates, and the nonempty chain's constant-singleton colimit is a singleton. The discrete two-object fold fails already on terminal objects, giving two points rather than one.
\end{proof}

Connectedness by itself is insufficient for product admission. A one-object category given by the nontrivial group $C_2$ has connected underlying category, so its fold preserves the terminal diagram. Colimit is the orbit functor. The regular $C_2$-set has one orbit, while the diagonal action on its Cartesian square has two orbits. Thus the product comparison is not bijective. This distinguishes the accumulation comparison from the earlier unit criterion.

\section{Precise reach}
This supplies a positive structure with genuine, possibly unbounded diagram shapes and some noninvertible shape changes. It does not supply arbitrary merges, arbitrary quotienting, or a runtime scheduling interpretation. The base extensions are projection-only; occurrence-changing content, shared effects and interference require additional interaction mechanisms. The theorem is an explicit family criterion and normal-form proof, not a claim that the full changing-shape contract has been met.
\paragraph{Literature and scope of attribution.} Siftedness and its finite-product characterization are established background; see Ad\'amek and Rosick\'y, \emph{On sifted colimits and generalized varieties}, Theory and Applications of Categories8(3),2001, and Ad\'amek, Rosick\'y and Vitale, \emph{What are sifted colimits?}, TAC23(13),2010. The natural-number and group-orbit calculations above are included directly rather than invoking an unstated general preservation theorem. Indexed monoidal coherence and Kan-extension composition are standard machinery. The project-specific content is their explicit retained module/crossing assembly and exact separation from richer interaction hypotheses; no novelty claim is made.
\end{document}
